\CharacterSheet{lover}{Amante}{

\shortdescription{Um guerreiro formidável, protetor da família.}

\ObjectivesHeading{Objetivos}
\begin{itemize}
  \item Descobrir quem matou \Character{newcomer} e conseguir que essa pessoa seja expulsa por votação. (\PositiveReward{1}, \NegativeReward{1})
\end{itemize}

\LosingConditionsHeading{Condições de derrota}
Não conseguir expulsar por votação quem matou \Character{newcomer}, caso haja uma pessoa responsável pelo homicídio.

\SecretsHeading{Segredos}
\begin{itemize}
  \item Mantenho uma relação secreta com \Character{newcomer} há um ano. (\NegativeReward{1}, \PositiveReward{1})
\end{itemize}

\ActionsHeading{Ações}
\begin{itemize}
  \item Sê direto, firme e fisicamente imponente. Sê o mais másculo dos homens.
  \item Protege a imagem de \Character{traitor}.
\end{itemize}

\BackgroundHeading{Contexto}

\timestamp{No ano 700 d.C.}, era um guerreiro que tinha passado a vida mortal a proteger quem não se conseguia defender. \Character{vampirehead} viu utilidade nessa força e acolheu-me na sua família como protetor. Aceitei porque a imortalidade me prometia os meios para continuar a lutar para sempre. Quando fui transformado, \Character{vampirehead}, \Character{usurper}, \Character{enforcer}, \Character{buttler}, \Character{ghoul} e \Character{supremacist} já faziam parte da família. A festa até foi bastante porreira.

\timestamp{No ano 1000 d.C.}, \Character{prevlover} juntou-se a nós como investigador da família.

\timestamp{No ano 1400 d.C.}, pedi a \Character{vampirehead} que admitisse \Character{traitor}. Amava-a e acreditava que o seu hábito de ajudar as pessoas comuns tornaria a nossa família menos cruel. Os outros viam nela uma humana caridosa que não tinha lugar entre vampiros, sobretudo \Character{supremacist}. Insisti até \Character{vampirehead} a aceitar.

Ao longo dos séculos, a preocupação de \Character{traitor} com os humanos tornou-se mais intensa e mais secreta. Tenho reparado em viagens sem explicação e em correspondência cuidadosamente escondida. Receio que esteja envolvida em algo perigoso. Não lhe perguntei diretamente porque ainda a amo e porque tenho medo da resposta.

\timestamp{Há dois anos}, \Character{vampirehead} começou a trazer \Character{newcomer} aos encontros de família. Conheci-o como o convidado predileto da chefe e, ao princípio, mantive a distância. Depois começámos a conversar. Há um ano, a amizade tornou-se uma relação secreta. Ele era humano, eu era vampiro, e as regras da família faziam disso motivo suficiente para nos arruinar aos dois, caso se soubesse.

\timestamp{Há dois meses}, \Character{vampirehead} anunciou que \Character{newcomer} se juntaria à família numa festa em \timestamp{16/7/2026}. Todos os vampiros tinham de comparecer.

\timestamp{Umas duas semanas antes da festa}, \Character{vampirehead} perguntou se mais alguém, além dela, de \Character{buttler} e de \Character{newcomer}, queria ajudar. Disse que \Character{usurper} se tinha oferecido e chegaria ao castelo duas noites antes da festa.

\timestamp{Desde as 16:00 de ontem}, tenho notado problemas de memória. Não é só ter-me esquecido de coisas; também me custa guardar novas recordações.

\timestamp{Ontem, às 21:00}, telefonei a \Character{newcomer} e pedi-lhe que se encontrasse comigo lá fora. Disse-me que a nossa relação tinha de acabar. Acreditava que, depois de \Character{vampirehead} o transformar, a chefe esperaria que os dois tivessem uma relação amorosa. Disse que talvez houvesse outra oportunidade para nós. Em vez de nos despedirmos, passámos uma última noite juntos.

\timestamp{Hoje de manhã, às 11:00}, implorei a \Character{newcomer} que não acabasse tudo de forma tão abrupta. Falámos lá fora, onde pensei que ninguém nos ouviria. Acho que alguém nos viu, mas os pormenores estão enevoados. Por alguma razão, não me consigo lembrar bem.

\GMNotes{
\paragraph{Porque se pode suspeitar desta personagem?}
\Character{lover} mantinha uma relação proibida com \Character{newcomer}, passou a noite com ele pouco antes da festa e pagou a \Character{buttler} para esconder uma conversa privada. A sua raiva e o seu desgosto tornam-no um suspeito emocionalmente credível. A ligação a \Character{traitor} também gera suspeitas quando a correspondência dela com a igreja é encontrada.

\paragraph{Como se pode provar a inocência desta personagem?}
O plano fatal já estava em marcha ao almoço do dia anterior, quando o vinho de \Character{usurper} deu secretamente a \Character{newcomer} a primeira dose da poção de transformação. \Character{lover} não participou nesse almoço nem mexeu no vinho. Os encontros com \Character{newcomer} aconteceram muito depois, às \timestamp{21:00} e na manhã seguinte, pelo que não explicam a primeira dose escondida.

\paragraph{Qual é o papel desta personagem na história e no jogo?}
\Character{lover} é um suspeito emotivo, mas inocente. O jogador tem de equilibrar a lealdade a \Character{newcomer} e a \Character{traitor} com os factos que incriminam ambas as relações. O pagamento dá a \Character{buttler} um motivo para admitir o que viu, enquanto a preocupação com \Character{traitor} dá um custo pessoal à conspiração da igreja.

\paragraph{Como pode esta personagem encaminhar a investigação para a solução?}
Se \Character{lover} admitir a conversa da manhã, o grupo pode descobrir que \Character{buttler} e \Character{usurper} estavam por perto. Interrogar \Character{buttler} sobre esse segredo pode levá-la a contar o que aconteceu ao almoço, quando \Character{usurper} substituiu o vinho preparado para \Character{newcomer}. Essa sequência chama a atenção para a garrafa anterior, em vez da relação de \Character{lover} com a vítima, que ocorreu mais tarde.
}
}
